\documentclass[11pt]{article}
\usepackage[margin=1in]{geometry}
\usepackage{amsmath,amssymb}
\usepackage{booktabs}
\usepackage{graphicx}
\usepackage{float}

\title{4/23: Eigenfunction Results}
\author{}
\date{}

\begin{document}
\maketitle

\section{Setup}

We compute the first nontrivial eigenfunction $Q_1(x)$ of the generator $L$ for the 3-mode stochastic NLS energy cascade:
\[
L Q_1 = \lambda_1 Q_1.
\]
Transition kernel:
\[
P(y, t \mid x, 0) \sim \sum_i e^{-\lambda_i t}\, Q_i(x)\, P_i(y),
\]
so for any observable $f$,
\[
\mathbb{E}_x[f(X_t)] = \sum_i \langle f, P_i \rangle\, Q_i(x)\, e^{-\lambda_i t}.
\]
At large $t$, the dominant nontrivial mode gives
\[
\mathbb{E}_x[f(X_t)] - \langle f \rangle_{\mathrm{ss}} \approx Q_1(x)\, e^{\lambda_1 t},
\]
where $\lambda_1 = \lambda_R + i\lambda_I$ with $\lambda_R < 0$.

\section{Observables}

We test two observables:
\[
f_A(x) = \max(0,\, 3 - I_1) + \max(0,\, 3 - I_2) + \max(0,\, 3 - I_3)
\]
\[
f_B(x) = \cos\theta_1 = \cos\bigl(2(\phi_1 - \phi_2)\bigr)
\]

\subsection{Single trajectory}
\begin{figure}
    \centering
    \includegraphics[width=1\linewidth]{4_23/trajectory_bump.png}
\end{figure}
\begin{figure}
    \centering
    \includegraphics[width=1\linewidth]{4_23/trajectory_cos.png}
\end{figure}
\section{MC Parameters}

\begin{center}
\begin{tabular}{lc}
\toprule
Parameter & Value \\
\midrule
$\gamma$ & 0.1 \\
$T_1 = T_3$ & 5.0 (equilibrium) \\
$I_{\max}$ & 10.0 \\
$\mathrm{d}t$ & 0.001 \\
Recording gap & 100 steps (= 0.1 time units) \\
$T_{\mathrm{steps}}$ & 50 (total time = 5.0) \\
$N_{\mathrm{init}}$ & 4,096 (5D LHS) \\
$N_{\mathrm{sample}}$ & 10,000  \\
EM steps per trajectory & 5,000 \\
\bottomrule
\end{tabular}
\end{center}

\section{Results}

\subsection{Eigenvalue}

\begin{center}
\begin{tabular}{lcc}
\toprule
 & Observable A (bump) & Observable B ($\cos\theta_1$) \\
\midrule
$\lambda_R$ & $-1.35 \pm 0.67$ & $-1.57 \pm 0.89$ \\
$\lambda_I$ & $\approx 0$ & $\approx 0$ \\
$\tau = 1/|\lambda_R|$ & 0.74 & 0.64 \\
Good fits & 3,743 / 4,096 (91\%) & 1,026 / 4,096 (25\%) \\
\bottomrule
\end{tabular}
\end{center}

\subsection{Check}
From the spectral expansion, after subtracting the steady state value, we have
\[
\mathbb{E}_x[f(X_t)] - \langle f \rangle_{\mathrm{ss}} \approx Q_1(x)\, e^{\lambda_1 t}.
\]

If $\lambda_1$ is real: $e^{\lambda_R t}$ is pure exponential decay. 

If $\lambda_1$ is complex: $e^{(\lambda_R + i\lambda_I)t}$ gives decayed oscillation. 

We check the 20 initial points with strongest signal. 

\begin{figure}[H]
    \centering
    \includegraphics[width=1\linewidth]{4_23/Figure_1.png}
\end{figure}



\end{document}